\chapter{泡沫与崩溃的通用剧本}
\chaptermeta{17}

\begin{ledetext}
四百年来它只演过一出戏，换的只是布景和主角的行头。剧本你看完就会背，难的是知道自己坐在第几幕。
\end{ledetext}

\section{牛顿的那笔账}

1720 年，艾萨克·牛顿在南海公司的股票上亏掉了大约两万英镑。当时他是英国皇家铸币厂厂长，年收入两千英镑上下。

二十年薪水。

流传最广的那句"我能计算天体的运行，却算不出人类的疯狂"，出处一直存疑，最早的记录出现在他死后很多年。但那笔亏损是有账可查的。

更要命的是细节：他不是从头套到尾的。1720 年初他买入，四月前后卖出，赚了大概七千英镑，本金翻了一倍。然后他站在场外，看着股价继续往上冲。六月他又进去了，这次买在接近顶部的位置。

让他破产的不是贪婪。是眼看着别人赚，而自己已经下车。

\section{五幕}

海曼·明斯基和查尔斯·金德尔伯格把这出戏拆成了五幕。这个结构从 1637 年用到 2026 年，一次没换过。

\begin{flowlist}
  \flowitem{01}{位移}{一件真事发生了。蒸汽机、铁路、电力、互联网、页岩油、大模型。它真的改变了利润的分布，也真的创造出一个以前不存在的赚钱机会。\textbf{泡沫的种子从来不是谎言，是真话。}}
  \flowitem{02}{繁荣}{信贷跟上来。银行、券商、基金开始往这个方向配钱。价格涨，早进去的人赚到了，而"有人真的赚到钱"这个事实，比任何论证都有说服力。}
  \flowitem{03}{狂热}{新人涌进来。估值和现金流脱钩，开始改用新指标：点击量、日活、GMV、算力储备。出现"这次不一样"的论述，而且这套论述通常有一半是对的，这才是它难对付的地方。}
  \flowitem{04}{获利了结}{最早进场的人开始减持。他们不发声明，他们只是卖。这个阶段价格往往还在涨，因为接的人比卖的人多。事后回看，这一幕最清楚；身在其中，这一幕最像"回调"。}
  \flowitem{05}{恐慌}{一件具体的、往往很小的事情触发它。流动性消失，杠杆被强平，强平压低价格，价格触发下一批强平。}
\end{flowlist}

\begin{chainflow}
\chainnode{价格下跌}\chainarrow \chainnode{追加保证金}\chainarrow \chainnode{被迫卖出}\chainarrow \chainnode{价格再跌}\chainarrow \chainnode{下一批被迫卖出}
\end{chainflow}

五幕之间没有钟点表。第三幕可以持续三个月，也可以持续四年。"识别泡沫"和"知道它什么时候破"是两个难度差着一个数量级的问题，绝大部分把这两件事混为一谈的人，都在第三幕里亏掉了钱。

\section{四个零件，缺一个就不是泡沫}

一个真实的新故事，便宜的信贷，杠杆，大量新参与者。四样凑齐才叫泡沫，少一样就只是一次普通的上涨。

第一样最反直觉，所以多说两句。

铁路是真的。1840 年代英国铁路狂热崩掉之后，英国确实得到了世界上第一张全国铁路网，那些铁轨后来跑了一百多年。互联网是真的，2000 年之后光纤、机房和一整代工程师留在了地上，被下一批公司用得很爽。AI 也是真的。

泡沫化的从来不是那个故事。

\begin{pullquote}{}
市场很少猜错方向，它经常猜错速度：把十五年才能兑现的现金流，按三年兑现折进今天的价格里。
\end{pullquote}

\begin{egbox}{把"时间表"算成数字}
一家公司今年收入 12 亿，市场给 40 倍市销率，市值 480 亿。

这个定价隐含的意思是：它得把收入做到 80 亿以上，并且做出 20\% 的净利率（16 亿利润），才能落到 30 倍市盈率这种"不算离谱"的位置。收入要涨到 6.7 倍。

每年增长 45\%，需要 5 年零 2 个月。\\ 每年增长 25\%，需要 8 年零 6 个月。

同一家公司，同一个故事，增速差 20 个百分点，兑现时间差三年多。而市场愿意等的时间，通常比它嘴上说的短得多。第三幕的崩塌几乎从来不是因为"故事被证伪"，是因为\textbf{某个季度的增速从 45\% 掉到了 25\%}。

\end{egbox}

\section{点名，每个只说最要命的那一句}

\textbf{1637 年，荷兰郁金香。}这个故事你一定听过：一颗球茎换一栋运河边的房子。

近二十年的档案研究给出的图景要冷淡得多。卷进去的主要是一个不大的、彼此认识的商人和手艺人圈子；成交大量是期货合约而不是现货交割；价格崩溃之后，荷兰经济没有出现可观察的衰退；法院基本以"这类合约不可强制执行"了事。它的经济影响被后世严重夸大了。

它被引用了四百年，是因为它简短、荒诞、好用。真正的教训不在球茎里，在你为什么会毫不怀疑地接受一个自己从没核对过来源的故事。这一段我特意放在最前面：怀疑现成叙事，本身就是这一章要教的东西。

\textbf{1720 年，南海公司。}牛顿那笔账。补一句：南海公司的商业模式是接管英国国债，它和政府的关系比它和南海贸易的关系密切得多。

\textbf{1929 年，美股。}要害是保证金。

\begin{egbox}{10\% 首付买股票}
1928 年，你拿出 1000 美元，按 10\% 的保证金买入 10000 美元的股票，剩下 9000 向券商借。

股票涨 20\%，市值 12000，还掉 9000，你手里剩 3000。本金翻两倍。\\ 股票跌 10\%，市值 9000，你的 1000 归零，追加保证金的电话打过来。\\ 股票跌 12\%，你不但归零，还倒欠券商 120。

1929 年 10 月底那几天，指数单日跌幅超过 12\%。电话打不通的结果是券商直接平仓，平仓压低价格，价格触发下一批平仓。第 14 章那台绞肉机的历史现场，就在这儿。

\end{egbox}

\textbf{1990 年，日本。}顶点那几年流传过一个说法：东京都的土地按市价算，卖掉可以买下整个美国。这个算法本身很粗糙，但它当时被广泛相信，这件事比算得准不准更重要。日经 225 在 1989 年 12 月 29 日收于 38957 点，此后花了 34 年才回到这个位置。

\textbf{2000 年，纳斯达克。}5048 点跌到 1114 点，跌掉 78\%，回到 5048 用了 15 年。

同一时期，亚马逊从最高点跌了约 94\%，从一百多美元跌到五块多。后来它涨了几百倍。

同一个泡沫，同一个方向判断（"互联网会改变世界"这句话完全正确），标的选择决定了两种完全相反的命运。而这两种命运对普通人的意义又完全不同：拿着一篮子纳斯达克，考验的是耐心；拿着亚马逊，考验的是你在跌掉 94\% 的那段路上会不会被迫卖出。

被迫卖出。这四个字才是重点。你有没有杠杆，会不会失业，家里有没有别的钱要用，决定了你能不能活到兑现那一天。\hlmark{在泡沫里，"看对趋势"的价值远低于"活到兑现"。}而后者取决于杠杆和现金流，不取决于眼光。

\textbf{2008 年，次贷。}一句话讲完：把一个统计观察（美国房价从未在全国范围内同时下跌）当成了物理定律，用它给证券定价，请评级机构盖章，然后在上面叠了二三十倍杠杆。假设错一次，杠杆把这个错误放大二三十倍还给你。

\textbf{2023 年，硅谷银行。}它没有坏账。它买的是美国国债和机构 MBS，信用风险接近于零。它错在久期：用随时可以提走的存款，买了平均久期五六年的债。利率从零升到 5\% 以上，债价跌，浮亏挂在"持有至到期"科目里不用计提，直到它必须卖出变现的那天。

真正的新东西是速度。2023 年 3 月 9 日一天被提走 420 亿美元，第二天预定还要走 1000 亿。1930 年代的挤兑要排队，2023 年的挤兑是几个微信群和 Slack 频道里的截图，加上手机银行的转出按钮。这是四百年里唯一一条被技术真正改写的规则。

\begin{tablewrap}
\begin{xltabular}{\linewidth}{>{\raggedright\arraybackslash}p{0.20\linewidth}XrX}
\toprule
\bthead{泡沫} & \bthead{顶到底} & \bthead{最大跌幅} & \bthead{回到原点用了} \\
\midrule
\textbf{1929 道琼斯} & 1929/09 – 1932/07 & -89\% & 25 年 \\
\textbf{1989 日经 225} & 1989/12 – 2009/03 & -82\% & 34 年（2024 年 2 月） \\
\textbf{2000 纳斯达克} & 2000/03 – 2002/10 & -78\% & 15 年 \\
\textbf{2007 标普 500} & 2007/10 – 2009/03 & -57\% & 5 年半 \\
\bottomrule
\end{xltabular}
\end{tablewrap}

看最后一列。这一列决定的不是收益率，是你的人生安排：25 年和 5 年半，对一个 55 岁准备退休的人来说是两个物种。

\bookbreak

\section{破了之后，只有三种结局}

\textbf{温和出清。}价格跌下来，杠杆低的机构接手便宜资产，两三年翻篇。1987 年的股灾是这一类，一天跌 22\%，第二年就修复了。前提是杠杆没有渗进银行体系。

\textbf{资产负债表衰退。}日本那一种。企业的资产跌破了负债，但现金流还在，于是它们把经营目标从"利润最大化"改成"负债最小化"：赚来的钱全部拿去还债，不再借钱扩张，哪怕利率是零。

这时候货币政策失效，不是因为央行不努力，是因为问题不在钱贵，在没人想借。辜朝明在 1990 年代提出这个描述时不被主流接受，现在它被广泛引用。日本用了差不多三十年才走出来。

\textbf{系统性危机加政府兜底。}2008 年那一种。政府接管、注资、担保，把私人部门的债务转成公共部门的债务。危机被中止了，代价是政府债务率永久性抬高一个台阶，而且回不去。美国联邦债务占 GDP 的比重在 2007 年是 60\% 出头，此后再没回到那个水平；2026 年 8 月，美国国债规模突破 40 万亿美元。

三种结局的分野只有一条：\textbf{杠杆有没有渗进银行体系。}没渗进去，是财富重新分配；渗进去了，是支付系统出问题。这条线在第 14 章画过一次，这里再画一遍。

\section{稳定本身孕育不稳定}

明斯基最反直觉的一句话，也是他最重要的一句。

\begin{deepbox}{进阶：可以跳过}
明斯基把借款人分成三类。\textbf{对冲型}：现金流能覆盖利息和本金。\textbf{投机型}：现金流只够付利息，本金到期要靠续借。\textbf{庞氏型}：现金流连利息都不够，全靠资产升值和再融资撑着。

他的论点是：一个经济体在长期平稳中，会自动从第一类往第三类漂移。不是因为人变坏了，是因为前面几年里谁激进谁赚钱，激进的机构份额变大，保守的基金经理被赎回或者被换掉。市场在用最有效率的方式，筛选出最脆弱的结构。

\end{deepbox}

日子越安稳，人越敢借；越敢借，系统越脆弱；越脆弱的系统在没出事之前，看起来越稳。

所以"这次很稳"这句话本身就是一个风险指标。听到它的频率越高，位置越靠后。

\section{AI 是不是泡沫}

这是 2026 年被问得最多的问题。我不会给答案，但我会把两边的证据都摆出来，最后给一个你自己能用的框架。

\begin{statrow}{3}
  \statitem{57}{\%}{2025 年数据中心资本开支同比增速}
  \statitem{2000}{亿美元}{私募信贷投向 AI 领域的规模}
  \statitem{78}{\%}{纳斯达克 2000–2002 最大跌幅（参照）}
\end{statrow}

\textbf{说它不是泡沫的证据。}2025 年数据中心资本开支同比增长约 57\%，出钱的主力是现金流极强的巨头，不是靠融资活着的初创公司。收入是真的，需求是真的，算力目前仍然供不应求。美联储主席凯文·沃什在 2026 年 7 月公开表示，人工智能相关投资正在为未来经济增长奠定基础。

\textbf{值得警惕的信号。}不止一条：

\begin{warnbox}{小心这四件事}
\textbf{折旧年限的会计假设。}把 GPU 的折旧年限从三年拉长到五年六年，当期利润立刻好看。如果这批芯片的实际有效寿命只有三年，多出来的那部分利润是从未来借的。

\textbf{供应链上的循环交易。}A 投资 B，B 拿这笔钱买 A 的产品，A 的收入增长了。分开看每一笔都合规，合并看是左手倒右手。这种结构在 2000 年的电信设备行业出现过。

\textbf{债务融资的比例在上升。}早期是现金流出钱，现在开始发债、做结构化融资、找私募信贷。上一章那 2000 亿美元就在这里。零件二和零件三正在到位。

\textbf{市场集中度极高。}指数涨幅高度依赖少数几家公司，这意味着"分散投资"提供的保护比你以为的薄。

\end{warnbox}

\begin{talkbox}
  \talkline{沃什，2026 年 7 月}{"人工智能相关投资正在为未来经济增长奠定基础。"}
  \talkline{同一场发布会}{"资本支出供应副作用的发生时间难以预测。"}
  \talkline{翻译}{方向他认，节奏他不认。这两句话不矛盾，它们合起来正好是本章的主题：故事可能是真的，时间表可能是错的。}
\end{talkbox}

给你三个可以自己查、自己算的指标，不需要任何内幕消息：

\textbf{一，这一轮扩张的钱从哪来？}经营现金流的占比在上升还是下降？债务融资的占比在往哪个方向走？这个比例的拐点比总量重要得多。

\textbf{二，收入跟得上资本开支吗？}把两条曲线的同比增速画在一张图上。差距是在收窄还是在扩大？扩大了几个季度？

\textbf{三，新进场的人有没有加杠杆？}散户是用闲钱买，还是用融资余额买？相关主题的杠杆 ETF 规模在涨吗？零件四到位没有，看这个。

这三个问题我现在都没有确定答案。\textbf{判断框架给你，答案时间会给。}

\begin{mythbox}{常见误解}
\mvfrow{误解}{"读完这一章我就能提前认出泡沫，然后在破之前跑掉。"}
\mvfrow{事实}{事后看曲线，一切清清楚楚；事前你手里只有一个概率分布。格林斯潘 1996 年 12 月说出"非理性繁荣"的时候，纳斯达克在 1300 点附近，此后它又涨了将近三倍，2000 年 3 月才见顶。他说对了性质，说早了三年多，而早三年的代价可能比说错更大。你真正能做的只有两件：控制杠杆；保证自己在下跌 50\% 的时候不会被迫卖出。这不是投资建议，这是关于"什么做得到、什么做不到"的说明。}
\end{mythbox}

\begin{takeaway}{本章一句话}
泡沫需要四个零件：一个真实的新故事、便宜的信贷、杠杆、大量新参与者。故事通常是真的，错的是价格和时间表。

\begin{itemize}
  \item 五幕：位移、繁荣、狂热、获利了结、恐慌。幕与幕之间没有时刻表。
  \item 破裂后只有三种结局，分野在于杠杆有没有渗进银行体系。
  \item 稳定孕育不稳定。"这次很稳"是风险指标，不是安慰。
  \item 2000 年纳斯达克回本用了 15 年，同期亚马逊跌 94\% 后涨了几百倍。看对方向不等于活到兑现。
\end{itemize}

\end{takeaway}

\begin{bridgebox}
\textbf{下一章}泡沫破了会痛，但它还会再来。因为它不是孤立事件，它挂在一个更大的钟摆上，而那个钟摆有好几个，长度还不一样。
\end{bridgebox}

